\subsection{Request Lifecycle and Holds}\label{feat:lifecycle}

\subsubsection{Description and Priority}
The three Allocation Classes resolve at different speeds: Exclusive and Metered lines resolve instantly, Matched lines wait for a human approval. A Request must therefore neither lose the room while the people are approved, nor charge for GPUs nobody can use because the people were refused. The lifecycle answers this with Required and Optional lines and with Holds bounded by the Objection Window: when a Request enters Pending, every satisfiable Exclusive and Metered line is placed on Hold, tentative and uncharged, while the Matched lines await approval. The Request is Allocated when every Required line is satisfied, and in that one step all Holds commit and all Charges and Earnings post. If a Required line is refused or the Objection Window lapses, every Hold releases and the Request goes to Waiting. The states and transitions are drawn in \figref{fig:states}. This feature realizes \FE{lifecycle} and underlies \UC{draft}, \UC{approve}, \UC{escalate}, and \UC{waitlist}. \priority{High.}

\subsubsection{Stimulus/Response Sequences}
\begin{itemize}
  \item \textbf{Stimulus:} A Request is submitted. \textbf{Response:} The system evaluates every line and moves the Request to Allocated, Pending, or Waiting.
  \item \textbf{Stimulus:} The last outstanding Required approval arrives. \textbf{Response:} Holds commit, Charges and Earnings post, everyone involved is notified, all in one transaction.
  \item \textbf{Stimulus:} A Required approval is refused, or the Objection Window lapses. \textbf{Response:} Every Hold and its held credits release; the Request goes to Waiting; on lapse the outstanding approvals escalate.
  \item \textbf{Stimulus:} Capacity a Waiting Request needs becomes free, or an escalated approval arrives. \textbf{Response:} The system checks whether every Required line can now be satisfied at once and, if so, allocates the Request.
\end{itemize}

\subsubsection{Functional Requirements}
\fr{Lifecycle.States}{A Request shall be in exactly one of the states Draft, Submitted, Pending, Allocated, Waiting, Withdrawn, or Refused, with exactly the transitions of \figref{fig:states}. Allocated, Withdrawn, and Refused are final except that Allocated returns to Waiting when a Required line is Bumped.}
\fr{Lifecycle.Line.States}{Each Request Line shall be in exactly one of the states Open, Awaiting approval, Approved, Held, Waitlisted, Allocated, Dropped, Released, Bumped, or Expired; Appendix~\ref{app:data} defines each.}
\fr{Lifecycle.Evaluate}{On submission the system shall, in one transaction, determine for every Exclusive and Metered line whether it is satisfiable now (\FR{Exclusive.Conflict}, \FR{Metered.Pool.Check}, \FR{Metered.Quota.Check}) and for every Matched line whether an approval is needed, and shall move the Request as follows: to Allocated if no line needs approval and every Required line is satisfiable; to Pending if any Required or Optional line needs approval; otherwise to Waiting.}
\fr{Lifecycle.Hold.Place}{When a Request enters Pending, the system shall place a Hold on every satisfiable Exclusive and Metered line. A Hold shall block other Requests exactly as an allocation does and shall post no Charge (\BR{hold-bounded}).}
\fr{Lifecycle.Hold.Credits}{While a line is on Hold, its locked Quote shall count as held credits of the Team, so that concurrent Pending Requests cannot together exceed the Balance (\BR{affordability}).}
\fr{Lifecycle.Hold.Lifetime}{A Hold shall end in exactly one of these ways: it commits when the Request becomes Allocated; it releases when the Request becomes Waiting, Refused, or Withdrawn; or it releases when the Objection Window lapses. No Hold shall outlive the Objection Window of its Request (\BR{hold-bounded}).}
\fr{Lifecycle.Hold.NoneAfterLapse}{After a Request's Objection Window has lapsed the system shall place no further Hold for that Request; it shall be allocated only under \FR{Lifecycle.Waiting.Reevaluate}.}
\fr{Lifecycle.Window.Start}{The Objection Window shall start when the Request first enters Pending, shall have the length set by the Capacity Planner (example default two working days), and shall not restart when a Candidate is replaced or the Request re-enters Pending (\BR{objection-window}).}
\fr{Lifecycle.Window.Lapse}{When the Objection Window lapses with a Required approval outstanding, the system shall release every Hold, move the Request to Waiting, and escalate the outstanding approvals (\FR{Matched.Escalate}), within five minutes of the lapse (\peref{jobs}).}
\fr{Lifecycle.Allocate}{When every Required line is satisfied, the system shall in one transaction commit every Hold into an allocation, post a Charge equal to the locked Quote for every allocated Exclusive and Metered line, form the Loans and post the Earnings for approved Matched lines, drop Optional lines under \FR{Lifecycle.Optional.Drop}, move the Request to Allocated, and notify the Manager (\BR{commit-atomic}).}
\fr{Lifecycle.Optional.Drop}{At allocation, an Optional line that is unsatisfiable, unapproved, or no longer affordable after the Required lines are charged shall be dropped, never waitlisted, and the Manager shall be told which lines were dropped and why (\BR{required-optional}).}
\fr{Lifecycle.Refuse.Line}{When a Required Matched approval is refused, the system shall release every Hold and its held credits, move the Request to Waiting, and ask the requesting Manager for a replacement (\FR{Matched.Refuse.Replace}).}
\fr{Lifecycle.Waiting.Repending}{When the Manager of a Waiting Request selects a replacement Candidate while the Objection Window is still running, the system shall re-evaluate the Exclusive and Metered lines, place Holds on those now satisfiable, waitlist the Required ones that are not, request the new approval, and move the Request to Pending.}
\fr{Lifecycle.Waiting.Reevaluate}{Whenever capacity that a Waiting Request's line waits for frees, or an approval on it arrives, the system shall check whether every Required line of that Request can be satisfied at that moment; if so it shall allocate the Request under \FR{Lifecycle.Allocate} at the locked Quotes; if not, the freed capacity shall be offered to the next waiting line and the Request's lines shall keep their waitlist positions.}
\fr{Lifecycle.Refused}{A Request shall become Refused when a Required Matched line cannot be satisfied under \FR{Matched.Refuse.None}; every Hold releases and nothing is charged.}
\fr{Lifecycle.Bump.Return}{When a Required line of an Allocated Request is Bumped, the Request shall return to Waiting with that line at the head of its waitlist and every other line unchanged (\BR{bumped-line}).}
\fr{Lifecycle.Withdraw}{The Manager shall be able to withdraw a Request from Draft, Pending, or Waiting at any time (\FR{Draft.Withdraw}).}
\fr{Lifecycle.Audit}{Every change of Request state or line state shall be recorded with the actor (a Member, or the system for timed events), the change, and the timestamp, and the record shall be visible to the Request's Manager and to the Capacity Planner (\BR{audit}).}
\fr{Lifecycle.Notify}{The system shall send an in-app notification to the affected Manager when a Request enters Pending, Allocated, Waiting, or Refused, when a Hold expires, when a line is Bumped or satisfied from the waitlist, and when a Request escalates; to a lending Manager when an approval is requested; and to the Capacity Planner when a Request escalates (\ciref{inapp}).}
